\documentclass[10pt,a4paper]{amsart}
\usepackage[margin=19mm]{geometry}
\usepackage{amsmath,amssymb,mathtools}
\usepackage[scr=rsfs,cal=euler]{mathalfa}
\usepackage{xfrac}
\usepackage[unicode,hidelinks,bookmarks=false]{hyperref}
\newcommand{\ie}{i.e.\ }
\newcommand{\cf}{cf.\ }
\newcommand{\resp}{respectively\ }
\newcommand{\Subsection}{\subsection}
\providecommand{\nobreakdash}{\nobreak\hbox{-}}
\setlength{\emergencystretch}{2em}
\multlinegap0pt
\usepackage{stmaryrd}
\usepackage{multirow}
\usepackage{dsfont}
\usepackage{rotating}
\usepackage{algorithm}
\usepackage{algpseudocode}
\renewcommand{\baselinestretch}{1.2}
\newcommand{\blue}[1]{#1}
\newcommand{\mgt}[1]{#1}
\newcommand{\red}[1]{{\color{red}#1}}
\newcommand{\teal}[1]{#1}
\newcommand{\ve}{{\bf e}}
\newcommand{\vh}{{\bf h}}
\newcommand{\va}{{\bf a}}
\newcommand{\vb}{{\bf b}}
\newcommand{\vc}{{\bf c}}
\newcommand{\vw}{{\bf w}}
\newcommand{\vx}{{\bf x}}
\newcommand{\vy}{{\bf y}}
\newcommand{\vu}{{\bf u}}
\newcommand{\vv}{{\bf v}}
\newcommand{\vz}{{\bf z}}
\newcommand{\qu}[1]{\tilde{#1}}
\newcommand{\vqu}[1]{\tilde{\mathbf{#1}}}
\newcommand{\dq}[1]{\hat{#1}}
\newcommand{\vdq}[1]{\hat{\mathbf{#1}}}
\newcommand{\dqm}[1]{\hat{#1}}
\newcommand{\dqset}{\hat{\mathbb Q}}
\newcommand{\udqset}{\hat{\mathbb U}}
\newcommand{\cset}{\mathbb C}
\newcommand{\dcset}{\hat{\mathbb C}}
\newcommand{\cgroup}{{\mathbb C}^{\times}}
\newcommand{\dcgroup}{\hat{\mathbb C}^{\times}}
\newcommand{\ucset}{\mathbb T}
\newcommand{\udcset}{\hat{\mathbb T}}
\newcommand{\supp}{\operatorname{supp}}
\newcommand{\SOSrank}{\text{SOS-rank}}
\DeclareMathOperator{\vecop}{vec}
\newcommand{\R}{\mathbb{R}}
\newcommand{\N}{\mathbb{N}}
\newcommand{\rank}{\operatorname{rank}}
\newcommand{\T}{\top}
\newcommand{\X}{\mathcal{X}}
\newcommand{\A}{\mathcal{A}}
\newcommand{\I}{\mathcal{I}}
\newcommand{\B}{\mathcal{B}}
\newcommand{\PP}{\mathcal{P}}
\newcommand{\C}{\mathcal{C}}
\newcommand{\D}{\mathcal{D}}
\newcommand{\LL}{\mathcal{L}}
\newcommand{\OO}{\mathcal{O}}
\newcommand{\0}{\mathbf{0}}
\newcommand{\1}{\mathbf{1}}
\newcommand{\dd}{\mathbf{d}}
\newcommand{\ii}{\mathbf{i}}
\newcommand{\jj}{\mathbf{j}}
\newcommand{\kk}{\mathbf{k}}
\newcommand{\vq}{\mathbf{q}}
\newcommand{\vg}{\mathbf{g}}
\newcommand{\pr}{\vec{r}}
\newcommand{\pc}{\vec{c}}
\newcommand{\ps}{\vec{s}}
\newcommand{\pt}{\vec{t}}
\newcommand{\pu}{\vec{u}}
\newcommand{\pv}{\vec{v}}
\newcommand{\pn}{\vec{n}}
\newcommand{\pp}{\vec{p}}
\newcommand{\pq}{\vec{q}}
\newcommand{\pl}{\vec{l}}
\newcommand{\vt}{\operatorname{vec}}
\newcommand{\sos}{\mathrm{sos}}
\newcommand{\bsr}{\mathrm{BSR}}
\newtheorem{Thm}{Theorem}[section]
\newtheorem{Def}[Thm]{Definition}
\newtheorem{Ass}[Thm]{Assumption}
\newtheorem{Lem}[Thm]{Lemma}
\newtheorem{conjecture}[Thm]{Conjecture}
\newtheorem{Prop}[Thm]{Proposition}
\newtheorem{Cor}[Thm]{Corollary}
\newtheorem{Conj}[Thm]{Conjecture}
\newtheorem{example}[Thm]{Example}
\newtheorem{remark}[Thm]{Remark}
\begin{document}
\title[Optimization Models and Computational Bounds for Limited Aug]{Optimization Models and Computational Bounds for Limited Augmented Zarankiewicz Numbers in the Incidence-Graph Family of Complete Graphs}
\author{Yi Xu, Gaohang Yu; Liqun Qi}
\date{}
\maketitle
\noindent\emph{Source and licence.} Optimization Models and Computational Bounds for Limited Augmented Zarankiewicz Numbers in the Incidence-Graph Family of Complete Graphs. Version arXiv:2605.29658v2 (CC BY 4.0). This material is distributed under the Creative Commons Attribution 4.0 International licence, http://creativecommons.org/licenses/by/4.0/, as recorded in the arXiv OAI licence metadata for this exact version and shown on the abstract page https://arxiv.org/abs/2605.29658v2. Full rights notice: licenses/arxiv-2605.29658-CC-BY-4.0.txt.\par\smallskip
\noindent\textit{Editorial scope.} Counted unit: the original \texttt{Prop} environment \texttt{-} at source lines 250--256 together with its complete proof at lines 258--282; the delivered document keeps source lines 210--283 so that every referenced label and every citation resolves inside it. Native editable source is the paper's own concatenated TeX; the AMS wrapper is editorial formatting only. No publisher class, upstream script or unrelated artwork is redistributed.
\section{Incidence-graph model and admissibility structure for \(K_{q+1}\)}

Fix an integer \(q\ge 2\). Let
\[
V=\{0,1,\dots,q\},\qquad S=\binom{V}{2},\qquad T=V.
\]
Define
\[
E_1=\{(e,v)\in S\times T: v\in e\}.
\]
Then
\[
G_1=(S,T,E_1)
\]
is the incidence graph of \(K_{q+1}\). It has
\[
\begin{aligned}
|S|&=\binom{q+1}{2},\qquad |T|=q+1,\\
|E_1|&=2\binom{q+1}{2}=q(q+1)=z\!\left(\binom{q+1}{2},q+1\right).
\end{aligned}
\]

\begin{example}
{The smallest nontrivial incidence-family case is \(q=2\). Then \(V=\{0,1,2\}\),}
\[
{S=\{01,02,12\},\qquad T=\{0,1,2\},}
\]
{and the incidence 1-edge set is}
\[
{E_1=\{(01,0),(01,1),(02,0),(02,2),(12,1),(12,2)\}.}
\]
{Thus each row in \(S\) has exactly one available cell, namely}
\[
{(01,2),\qquad (02,1),\qquad (12,0).}
\]
{This tiny case is useful only as a notation guide: it shows concretely what we mean by rows, columns, occupied cells, and available cells, and it also illustrates that a 2-edge is simply an unordered pair of available cells. The first genuinely nontrivial optimization phenomena occur for larger values of \(q\), but this \(q=2\) picture serves as a convenient model for the general definitions below.}
\end{example}

\blue{The following proposition shows that, for these parameters, the extremal \(C_4\)-free bipartite graph is unique up to isomorphism. Consequently, the choice of the 1-edge subgraph is unambiguous in the incidence-graph family under study.}


\begin{Prop}
\blue{Let \(n=q+1\) and \(m=\binom{n}{2}\). If \(H=(X,Y,F)\) is a \(C_4\)-free bipartite graph with}
\[
\blue{|X|=m,\qquad |Y|=n,\qquad |F|=n(n-1),}
\]
\blue{then \(H\) is isomorphic to the incidence graph of \(K_n\).}
\end{Prop}

\begin{proof}
\blue{For each \(x\in X\), let \(d(x)\) denote its degree. Because \(H\) is \(C_4\)-free, any two vertices of \(Y\) have at most one common neighbor in \(X\). Hence}
\[
\blue{\sum_{x\in X}\binom{d(x)}{2}\le \binom{n}{2}=m.}
\]
\blue{On the other hand, for every integer \(d\ge 0\),}
\[
\blue{\binom{d}{2}\ge d-1,}
\]
\blue{with equality if and only if \(d\in\{1,2\}\). Therefore}
\[
\blue{m\ge \sum_{x\in X}\binom{d(x)}{2}\ge \sum_{x\in X}(d(x)-1)=\sum_{x\in X}d(x)-|X|.}
\]
\blue{Now \(\sum_{x\in X}d(x)=|F|=n(n-1)\), while \(|X|=m=\binom{n}{2}=n(n-1)/2\). Thus}
\[
\blue{\sum_{x\in X}d(x)-|X|=n(n-1)-\binom{n}{2}=m.}
\]
\blue{So equality holds throughout. In particular, \(\binom{d(x)}{2}=d(x)-1\) for every \(x\in X\), hence \(d(x)\in\{1,2\}\) for all \(x\). Since the average degree on \(X\) is}
\[
\blue{\frac{|F|}{|X|}=\frac{n(n-1)}{\binom{n}{2}}=2,}
\]
\blue{every vertex of \(X\) must actually have degree exactly \(2\).}

\blue{Each vertex \(x\in X\) therefore determines an unordered pair of neighbors in \(Y\). Two distinct vertices of \(X\) cannot determine the same pair, for otherwise those two vertices together with those two vertices of \(Y\) would form a \(C_4\). Hence the map \(x\mapsto N(x)\) is an injection from \(X\) into \(\binom{Y}{2}\). But both sets have size \(m=\binom{n}{2}\), so this injection is a bijection. Thus every two-element subset of \(Y\) occurs exactly once as the neighborhood of a vertex of \(X\). This is precisely the incidence graph of \(K_n\).}
\end{proof}


\end{document}
